\documentclass[11pt]{article}
\newcommand{\DocNumber}{VOLUME 01}
\newcommand{\DocTitle}{C++23 Language Technical Reference}
\newcommand{\DocSubtitle}{Core-language semantics, clauses 1--15}
\newcommand{\DocSummary}{A concept-first technical reference with stable-name locators, diagnostic categories, and original code examples.}
% Shared, pdfLaTeX-only design system for the C++23 documentation suite.
\usepackage[T1]{fontenc}
\usepackage[utf8]{inputenc}
\usepackage{lmodern}
\usepackage[english]{babel}
\usepackage[letterpaper,margin=0.82in,headheight=16pt,headsep=17pt]{geometry}
\usepackage{microtype}
\usepackage{amsmath,amssymb}
\usepackage{array,booktabs,tabularx,longtable}
\usepackage[table]{xcolor}
\usepackage[most]{tcolorbox}
\usepackage{listings}
\usepackage{fancyhdr}
\usepackage{titlesec}
\usepackage{enumitem}
\usepackage{hyperref}
\usepackage{needspace}
\usepackage{etoolbox}
\definecolor{ink}{HTML}{13253B}
\definecolor{teal}{HTML}{006E7F}
\definecolor{muted}{HTML}{5A6C7C}
\definecolor{fog}{HTML}{EDF4F6}
\definecolor{sand}{HTML}{FAF3E8}
\definecolor{rule}{HTML}{C9D7DF}
\definecolor{codebg}{HTML}{F7F9FB}
\definecolor{codecomment}{HTML}{47715D}
\hypersetup{colorlinks=true,linkcolor=teal,urlcolor=teal,pdfauthor={Independent C++23 engineering documentation},pdftitle={C++23 Technical Suite}}
\setlength{\parindent}{0pt}
\setlength{\parskip}{6pt}
\setcounter{tocdepth}{2}
\renewcommand{\arraystretch}{1.18}
\titleformat{\section}{\Large\bfseries\color{ink}}{\thesection}{0.7em}{}
\titleformat{\subsection}{\large\bfseries\color{teal}}{\thesubsection}{0.55em}{}
\titlespacing*{\section}{0pt}{17pt}{7pt}
\titlespacing*{\subsection}{0pt}{11pt}{4pt}
\pagestyle{fancy}
\fancyhf{}
\fancyhead[L]{\footnotesize\bfseries\color{teal}C++23 / ISO 14882:2024}
\fancyhead[R]{\footnotesize\color{muted}\DocNumber}
\fancyfoot[L]{\scriptsize\color{muted}Independent engineering companion, not an ISO publication}
\fancyfoot[R]{\scriptsize\color{muted}\thepage}
\renewcommand{\headrulewidth}{0.35pt}
\lstdefinestyle{cppstyle}{language=C++,basicstyle=\ttfamily\footnotesize,keywordstyle=\color{teal}\bfseries,commentstyle=\color{codecomment},stringstyle=\color{ink},backgroundcolor=\color{codebg},frame=single,rulecolor=\color{rule},breaklines=true,breakatwhitespace=true,columns=fullflexible,keepspaces=true,showstringspaces=false,tabsize=2,numbers=none,upquote=true,xleftmargin=3pt,xrightmargin=3pt,aboveskip=6pt,belowskip=8pt}
\lstset{style=cppstyle}
\newtcolorbox{standardbox}{enhanced,breakable,colback=fog,colframe=teal!35!white,boxrule=0.55pt,arc=1.4mm,left=3mm,right=3mm,top=2mm,bottom=2mm,before skip=7pt,after skip=8pt}
\newtcolorbox{practicebox}{enhanced,breakable,colback=sand,colframe=ink!20!white,boxrule=0.55pt,arc=1.4mm,left=3mm,right=3mm,top=2mm,bottom=2mm,before skip=7pt,after skip=8pt}
\newcommand{\anchorref}[1]{\par\smallskip{\footnotesize\textbf{Standard locator:} \texttt{#1}}\par}
\newcommand{\standard}[1]{\begin{standardbox}\textbf{Standard-grounded interpretation.} #1\end{standardbox}}
\newcommand{\practice}[1]{\begin{practicebox}\textbf{Engineering recommendation (not a normative ISO rule).} #1\end{practicebox}}
\newcommand{\takeaway}[1]{\par\textbf{Review question.} #1\par}
\newcommand{\docfront}{%
\begin{titlepage}
\vspace*{25mm}
{\color{teal}\bfseries\large ENGINEERING REFERENCE SERIES\par}
\vspace{12mm}
{\raggedright\color{ink}\bfseries\fontsize{24}{30}\selectfont\DocTitle\par}
\vspace{7mm}
{\color{muted}\fontsize{15}{19}\selectfont\DocSubtitle\par}
\vspace{12mm}
\noindent\color{teal}\rule{\linewidth}{2pt}\par
\vspace{9mm}
{\large\DocSummary\par}
\vfill
\begin{tabularx}{\linewidth}{@{}lX@{}}
\textbf{Volume} & \DocNumber\\
\textbf{Standards basis} & ISO/IEC 14882:2024 (seventh edition, C++23)\\
\textbf{Document type} & Original interpretive and instructional material\\
\textbf{Date} & 16 October 2026\\
\textbf{Toolchain} & pdfLaTeX; illustrative C++23 examples require separate compiler support\\
\end{tabularx}
\vspace{8mm}
{\footnotesize This document is not an authorized copy of, replacement for, or amendment to ISO/IEC 14882:2024. The original standard controls in the event of disagreement.}\par
\end{titlepage}
\pagenumbering{roman}
\tableofcontents
\clearpage
\pagenumbering{arabic}
\section*{How to use this volume}\addcontentsline{toc}{section}{How to use this volume}
The reference is \emph{ISO/IEC 14882:2024, Programming languages --- C++}. Section locators and square-bracketed stable names identify where to verify precise wording; they are not quotations. \textbf{Standard-grounded interpretation} boxes summarize language or library semantics. \textbf{Engineering recommendation} boxes contain optional design, security, portability, or process guidance. Code is original and illustrative; individual library facilities may not yet be implemented by every vendor. This suite does not grant rights to reproduce or distribute the licensed standard.\par\medskip
}
\newcommand{\docend}{\clearpage\section*{Reference and traceability}\addcontentsline{toc}{section}{Reference and traceability}
\textbf{Primary source.} ISO/IEC 14882:2024(en), \emph{Programming languages --- C++}, seventh edition, October 2024. Use its numbered clauses, stable names, Annex C (informative compatibility), and Annex D (compatibility features) for authoritative details.\par
\textbf{Limits.} Clause-level summaries and exercises in this document are selective explanations, not exhaustive reproductions of all semantic constraints, library mandates, complexity clauses, and overload participation requirements. A successful compiler test demonstrates behavior of that specific implementation and configuration, not universal conformance.\par
\textbf{Use of contemporary engineering tools.} Compilers, sanitizers, analyzers, CI systems, repository policies, and vulnerability taxonomies are implementation guidance rather than language-standard requirements.\par}

\begin{document}
\docfront
\section{Reader contract and semantic status}
This is a clause-indexed interpretive reference for the core language, not a substitute for exact grammar, diagnostic obligations, or implementation documentation. Terms such as well-formed, ill-formed, undefined behavior, unspecified behavior, and implementation-defined behavior must not be conflated.

\small
\begin{longtable}{@{}p{.18\linewidth}p{.40\linewidth}p{.34\linewidth}@{}}
\toprule
\textbf{Classification} & \textbf{Meaning for analysis} & \textbf{Verification route}\\
\midrule
\endfirsthead
\toprule
\textbf{Classification} & \textbf{Meaning for analysis} & \textbf{Verification route}\\
\midrule
\endhead
Ill-formed, diagnosable & A conforming implementation must issue at least one diagnostic for a diagnosable violation. & 4.1.1; compiler diagnostics\\
\addlinespace[2pt]
Ill-formed, no diagnostic required & No reliable compiler rejection obligation; avoid constructions dependent on such ambiguities. & 4.1.1; relevant clause\\
\addlinespace[2pt]
Undefined behavior & The standard imposes no behavioral requirements for the affected execution. & Clause 3; 6.9; specific operation\\
\addlinespace[2pt]
Unspecified behavior & One of the allowed possibilities may occur, and which one need not be documented. & Clause 3; specific operation\\
\addlinespace[2pt]
Implementation-defined behavior & Choice must be documented by the implementation. & 4.1.2; implementation manual\\
\addlinespace[2pt]
Conditionally-supported & An implementation may elect not to support the construct; absence must be documented. & 4.1.1\\
\addlinespace[2pt]
\bottomrule
\end{longtable}
\normalsize
\section{Translation units, preprocessing and program structure}
\anchorref{4; 5.1--5.4; 15}
One source file is not the whole program. Translation phases, preprocessing, module/header importation, linkage, and the one-definition rule collectively determine program meaning.

\standard{The standard defines translation, linkage and diagnostic constraints independently of any particular compiler invocation. Header names and preprocessing tokens do not map one-to-one onto a build system.}
\begin{itemize}[leftmargin=*,itemsep=1pt,topsep=3pt]
\item Model translation units explicitly; document macro-visible interfaces.
\item Prefer intentional header or module interfaces over broad transitive inclusion.
\item Review ODR-sensitive inline definitions across configurations.
\end{itemize}
\begin{lstlisting}[language=C++]
// api.hpp
#pragma once
int normalize(int value);
// api.cpp
#include "api.hpp"
int normalize(int value) { return value < 0 ? 0 : value; }
\end{lstlisting}
\practice{Treat build reproducibility, generated headers and feature flags as part of the semantic input.}
\takeaway{How could two compilation configurations silently violate the one-definition rule?}
\section{Names, scope, lookup and linkage}
\anchorref{6.2--6.6 [basic.def] [basic.scope] [basic.lookup]}
A declaration introduces or redeclares entities; lookup determines which declarations are visible, and linkage determines whether names can denote the same entity across translation units.

\standard{Lookup, scope, overload sets, and the one-definition rule interact. Argument-dependent lookup may consider associated namespaces and classes beyond ordinary unqualified lookup.}
\begin{itemize}[leftmargin=*,itemsep=1pt,topsep=3pt]
\item Use explicit namespaces to represent ownership.
\item Avoid accidental ADL via unconstrained generic operator overloads.
\item Constrain public declarations and test multiple translation units.
\end{itemize}
\begin{lstlisting}[language=C++]
namespace domain {
struct Key { int value; };
constexpr bool operator==(Key a, Key b) { return a.value == b.value; }
}
static_assert(domain::Key{3} == domain::Key{3});
\end{lstlisting}
\takeaway{When does unqualified lookup differ from argument-dependent lookup?}
\section{Object model, storage duration and lifetime}
\anchorref{6.7 [basic.memobj] [basic.life]}
An object has a type, storage, and a lifetime. Acquiring storage alone does not license arbitrary typed access; lifetime start and end rules, union active members, and transparent replacement require care.

\standard{Lifetime, pointer validity, dynamic storage duration, and object representation are separate normative concepts. Access outside a permitted lifetime can have undefined behavior.}
\begin{itemize}[leftmargin=*,itemsep=1pt,topsep=3pt]
\item Make ownership and borrowing explicit at API boundaries.
\item Prefer automatic lifetime and standard containers over untyped storage.
\item Use special lifetime-management facilities only where their preconditions are met.
\end{itemize}
\begin{lstlisting}[language=C++]
#include <memory>
struct Record { int id; };
std::unique_ptr<Record> make_record(int id) {
    return std::make_unique<Record>(Record{id});
}
\end{lstlisting}
\takeaway{Why is a pointer value alone insufficient evidence that its referent is alive?}
\section{Types, cv qualification, pointers and references}
\anchorref{6.8; 7.3; 9.3}
Fundamental and compound types participate in conversions and overload resolution. A reference is not a general-purpose nullable owner, and top-level const differs from const of a referenced object.

\standard{Type identity, cv qualification, reference binding, and conversions are governed by distinct rules; volatile is not a threading synchronization primitive.}
\begin{itemize}[leftmargin=*,itemsep=1pt,topsep=3pt]
\item Use const to express non-mutation of the observable interface.
\item Use a pointer when absence is allowed; use references for required objects.
\item Make signed/unsigned crossings reviewable.
\end{itemize}
\begin{lstlisting}[language=C++]
#include <cstddef>
void inspect(const int& value);
void maybe_inspect(const int* value) {
    if (value != nullptr) inspect(*value);
}
\end{lstlisting}
\takeaway{Can a const-qualified pointer modify the object through a different alias?}
\section{Expressions, value categories and sequencing}
\anchorref{7.1--7.6; 6.9}
Expressions have types and value categories; conversions, temporary materialization, overloads and sequencing jointly determine evaluations and lifetime.

\standard{Not all operand evaluations are sequenced left-to-right. Separate indeterminately sequenced operations from unsequenced conflicting modifications; side-effect freedom is often the simplest invariant.}
\begin{itemize}[leftmargin=*,itemsep=1pt,topsep=3pt]
\item Do not rely on argument evaluation order without a guarantee.
\item Use named intermediates at boundary checks and state mutations.
\item Audit overloaded operators for hidden side effects.
\end{itemize}
\begin{lstlisting}[language=C++]
#include <vector>
int next_value();
void consume(int, int);
void safe_call() {
    const int first = next_value();
    const int second = next_value();
    consume(first, second); // explicit call order
}
\end{lstlisting}
\takeaway{What exactly does the comma operator guarantee that a comma separating arguments does not?}
\section{Constant evaluation and compile-time constraints}
\anchorref{7.7; 9.2.6 [dcl.constexpr]}
Constant expressions can reject invalid computations during translation; constexpr, consteval, and constinit have different roles.

\standard{Constant-evaluation rules restrict permissible operations and observable outcomes; constinit requires static/thread storage initialization to be static but does not imply constness.}
\begin{itemize}[leftmargin=*,itemsep=1pt,topsep=3pt]
\item Use static\_assert for numeric and type invariants.
\item Keep consteval functions small and diagnostic-oriented.
\item Reserve constinit for initialization-order guarantees where appropriate.
\end{itemize}
\begin{lstlisting}[language=C++]
consteval int square(int x) { return x * x; }
constexpr int capacity = square(8);
static_assert(capacity == 64);
constinit int global_count = 0;
\end{lstlisting}
\takeaway{Why is a constexpr variable not interchangeable with a constinit variable?}
\section{Declarations, initialization and structured bindings}
\anchorref{9 [dcl] [dcl.init]}
Initialization syntax can select different constructors and narrow conversions differently. Structured bindings introduce names whose reference and lifetime rules depend on their initializer.

\standard{Copy-, direct-, and list-initialization are not semantically interchangeable; narrowing restrictions apply to list-initialization.}
\begin{itemize}[leftmargin=*,itemsep=1pt,topsep=3pt]
\item Prefer braces when their initializer-list semantics are intended.
\item Audit implicit conversion constructors and overloaded initializer-list constructors.
\item Use explicit references when mutation of a bound object matters.
\end{itemize}
\begin{lstlisting}[language=C++]
#include <utility>
struct Bounds { int lo; int hi; };
constexpr Bounds b{2, 9};
static_assert(b.lo < b.hi);
auto [first, second] = std::pair{1, 2};
\end{lstlisting}
\takeaway{When does using braces choose a different overload than parentheses?}
\section{Functions, overload resolution and conversions}
\anchorref{9.5; 12}
Function viability, conversions, constraints, templates, and partial ordering drive call resolution. Default arguments and virtual dispatch do not obey the same selection model.

\standard{Overload resolution defines candidate/viable/best functions. A change in conversion ranking can change a previously valid call even when all overloads remain available.}
\begin{itemize}[leftmargin=*,itemsep=1pt,topsep=3pt]
\item Test overload resolution using static assertions and compile-negative cases.
\item Avoid large implicit-conversion surfaces in public APIs.
\item Define exception guarantees at function boundaries.
\end{itemize}
\begin{lstlisting}[language=C++]
void emit(int);
void emit(double);
void caller() {
    emit(12);    // exact match
    emit(12.0);  // exact match
}
\end{lstlisting}
\takeaway{Why can adding a new overload make a previously well-formed call ambiguous?}
\section{Class invariants, construction and inheritance}
\anchorref{11}
Classes support data encapsulation, special member functions, inheritance, access control and object initialization. Prefer invariants rather than merely grouping mutable members.

\standard{Base/derived construction and destruction have specified ordering; virtual dispatch during construction or destruction has phase-specific behavior.}
\begin{itemize}[leftmargin=*,itemsep=1pt,topsep=3pt]
\item Represent ownership with members that already express it.
\item Use override and final where they clarify interfaces.
\item Maintain invariants across construction, move, assignment and destruction.
\end{itemize}
\begin{lstlisting}[language=C++]
#include <vector>
class Buffer {
    std::vector<unsigned char> bytes_;
public:
    explicit Buffer(std::size_t n) : bytes_(n) {}
    std::size_t size() const noexcept { return bytes_.size(); }
};
\end{lstlisting}
\takeaway{What invariant must hold after a move operation on your class?}
\section{Templates, constraints and two-phase lookup}
\anchorref{13; 18 [temp] [temp.constr]}
Templates are declarations of families of entities. Deduction, instantiation, name resolution and constraints must be analyzed separately.

\standard{A satisfied syntactic concept is not necessarily modeled if semantic requirements are violated. Some template errors are ill-formed with no diagnostic required.}
\begin{itemize}[leftmargin=*,itemsep=1pt,topsep=3pt]
\item Name public concepts after semantic capabilities rather than syntax alone.
\item Test representative types, including intentionally invalid ones.
\item Inspect dependency of names used in templates.
\end{itemize}
\begin{lstlisting}[language=C++]
#include <concepts>
template<std::integral T>
constexpr T clamp_nonnegative(T x) { return x < 0 ? T{0} : x; }
static_assert(clamp_nonnegative(-4) == 0);
\end{lstlisting}
\takeaway{What does a requires-expression check, and what does it not prove?}
\section{Modules and interface reachability}
\anchorref{10 [module]}
Module units, named modules, imports, exports, global module fragments and reachability establish interfaces beyond the conventional header model.

\standard{The grammar and visibility rules are normative; compiler build ordering, BMI artifacts and command lines are implementation-specific.}
\begin{itemize}[leftmargin=*,itemsep=1pt,topsep=3pt]
\item Establish a dependency graph and build exported interfaces first.
\item Document macro-to-module migration concerns.
\item Use feature-gated module trials until the selected vendor is confirmed.
\end{itemize}
\begin{lstlisting}[language=C++]
// math.ixx (illustrative module interface)
export module math;
export int double_value(int x) { return 2 * x; }
// consumer.cpp
// import math;
// int main() { return double_value(5) != 10; }
\end{lstlisting}
\takeaway{Why can exporting a declaration differ from making it reachable?}
\section{Lambdas, explicit object parameters and callable objects}
\anchorref{7.5.5; 9.3.4.6}
C++23 supports an explicit object parameter for suitable non-static member functions, allowing generic forms of the object argument. Lambdas capture by value/reference subject to lifetime and mutability constraints.

\standard{Capture initialization, closure lifetime, and explicit-object syntax are governed by the core language; an explicit object parameter is not an implicit this pointer.}
\begin{itemize}[leftmargin=*,itemsep=1pt,topsep=3pt]
\item Prefer value capture for detached asynchronous tasks, after considering object cost.
\item Avoid capturing temporaries by reference when callbacks may outlive them.
\item Use explicit-object functions when they simplify cv/ref overload duplication.
\end{itemize}
\begin{lstlisting}[language=C++]
struct Counter {
    int value{};
    constexpr int read(this const Counter& self) { return self.value; }
};
static_assert(Counter{7}.read() == 7);
\end{lstlisting}
\takeaway{Which compiler feature test should you evaluate before deploying this syntax?}
\section{Exceptions, unwinding and noexcept}
\anchorref{14}
Exception propagation interacts with destructors, construction, exception specifications and the termination model.

\standard{An exception leaving a noexcept(true) function causes termination. Exception safety is a library/client design property, not a guarantee inferred from noexcept on every operation.}
\begin{itemize}[leftmargin=*,itemsep=1pt,topsep=3pt]
\item Use RAII for cleanup on every exit path.
\item Specify basic or strong exception guarantees where needed.
\item Avoid exceptions escaping destructors during unwinding.
\end{itemize}
\begin{lstlisting}[language=C++]
#include <mutex>
void update(std::mutex& mx, int& shared) {
    std::lock_guard<std::mutex> lk{mx};
    ++shared; // lock released on normal and exceptional exit
}
\end{lstlisting}
\takeaway{Why is a noexcept annotation on a function not equivalent to a proof of exception safety?}
\section{Coroutines and suspension-point ownership}
\anchorref{9.5.4; 17.12; 26.8}
Coroutines transform control flow around promise objects, coroutine state and suspension. Library generators are one consumer of this core facility.

\standard{The coroutine machinery and awaitable protocol are specified, but the safety of references stored across suspension points depends on user code and lifetime.}
\begin{itemize}[leftmargin=*,itemsep=1pt,topsep=3pt]
\item Treat each suspend point as a potential reentrancy and lifetime boundary.
\item Avoid retaining borrowed stack references across independently scheduled operations.
\item Test cancellation/destruction at every suspension state.
\end{itemize}
\begin{lstlisting}[language=C++]
// Coroutine lifetime checklist (pseudocode):
// 1. Who owns the coroutine handle/frame?
// 2. Does each referenced object outlive every suspension?
// 3. What destroys the frame on success, error, and cancellation?
\end{lstlisting}
\takeaway{What state survives suspension and who releases it?}
\section{Core-language review checklist}
\begin{itemize}[leftmargin=*,itemsep=1pt,topsep=3pt]
\item Identify lifetime owner, storage duration and aliasing assumptions for each exported API.
\item Trace each template error to substitution, viability or a hard diagnostic.
\item Review the same program across multiple translation units and optimization modes.
\item Record which properties are guaranteed, implementation-defined, unspecified or outside the standard.
\end{itemize}
\docend
\end{document}
